% No preamble. Prefix eqr-. Infinite local quotient; finite-history boundary explicit.
\section{Reference extraction and a contractive finite-block recurrence}
\label{eqr-section}

Work on the selected joint homogeneous classical system of
Sections~\ref{csgq-section} and \ref{qll-section}. Fix the primitive
odd ratio $\ell$, the elementary composite ratio $R=\ell^s$, the
auxiliary grouping $J$, and carrier exponent $a_0=1/16$.
At stage $j$ the preceding depth is $M_j=M_0R^j$. Each stage uses
the same cell grouping and analytic radius in its own current units.
All constants below may depend on this fixed geometry, but are uniform
in $M_j$. Use a simple gauge Lie algebra, in particular
$\mathfrak{su}(2)$, so the Maxwell functional is scalar. The finite
family of componentwise references gives the corresponding version
for a semisimple product.

The local quantum quotient map is an established input from
Theorem~\ref{qll-descent-theorem}, with its strict domain, periodization,
classical-system and logarithm-branch hypotheses. The discussion below
does not replace that theorem by an inference from compact-field
equality alone. The actual finite-regulator history remains a separate
issue stated at the end.

\subsection{The actual reference transports exactly}
\label{eqr-reference-section}

Let $U_M(D)$ be the selected ordinary homogeneous minimum and let
\begin{equation}
 S_M(D)=S_{\rm bare}(U_M(D)),\qquad S_M(1)=0,
 \qquad
 P_{M,R}(W)=\mathcal M_\ell^k U_{MR}(W),\quad M=\ell^k.
 \label{eqr-reference-definition}
\end{equation}
The trace convention and plaquette counting are fixed, and the coupling
is outside $S_M$. Classical background substitution is
\[
        (\mathcal T_{M,R}E)(W)=E(P_{M,R}(W)).
\]
It contains no fluctuation integration.

\begin{lemma}[Exact reference identity]
\label{eqr-reference-lemma}
On the specified common local analytic domains in dimension four,
\begin{equation}
                  \mathcal T_{M,R}S_M=S_{MR}.
 \label{eqr-reference-transport}
\end{equation}
This identity holds for the actual ordinary classical reference,
not merely for a reference with the same quadratic normalization.
It passes to the jointly chosen infinite normalized classes.
\end{lemma}
\begin{proof}
On a finite compatible regulator put
\[
 U_2=U_{MR}(W),\qquad D=\mathcal M_\ell^kU_2,\qquad U_1=U_M(D).
\]
Both fine configurations are fields on the same oriented graph.
Its fine spacing is $1/M$ in $D$ units and $1/(MR)$ in $W$ units.
Exact ordinary composition gives
\begin{equation}
 \mathcal M_\ell^kU_1=D=\mathcal M_\ell^kU_2,\qquad
 \mathcal M_\ell^{k+s}U_1=W=\mathcal M_\ell^{k+s}U_2.
 \label{eqr-two-constraints}
\end{equation}
The actual minimum comparison balls can be made common for fixed $R$.
Indeed a fine physical potential changes units by
$A_1^W(x)=R A_1(Rx)$, so
\begin{equation}
 \|A_1^W\|_{C^{1,\alpha}}
       \le CR^{2+\alpha}\|A_1\|_{C^{1,\alpha}}.
 \label{eqr-reference-rescaling}
\end{equation}
The prefix bound is
$\|\log D/i\|_\infty\le C_A\|\log W/i\|_\infty/R$.
Thus, for fixed $R$, a sufficiently small real $W$ ball places
both fields in both needed comparison classes, uniformly in
preceding refinement and volume.

The minimum at $D$ gives $S_{\rm bare}(U_1)\le S_{\rm bare}(U_2)$;
the minimum at $W$ gives the reverse inequality. In dimension four
the coupling-free bare action is the same plaquette sum under this
coordinate relabeling: its physical prefactor is
$\eta^{d-4}=1$. Hence the two actions agree. This is also the identity
$S_M(V^\star)=S_M(V_p)=S_{MR}$ in the actual common-domain
factorization.

Both sides are holomorphic on the specified connected common domain.
Uniqueness from the real germ continues the identity there.
For $R>C_A$, the angle estimate maps the common $b$ ball into itself;
the smaller real comparison ball does not cause a repeated analytic
radius loss. No minimum assertion at complex fields is used.

For the fixed joint infinite system, pass vacuum-normalized values on
compact $W$ using the connected interaction bounds and summable-field
continuity. These are the same hypotheses used for the classical
substitution limit. No convergence of an extensive infinite action
value is asserted. The finite references vanish at the trivial field,
so the normalized identity has no vacuum ambiguity.
\end{proof}

The exact flat midpoint symbol and its identification with the
classical Hessian give
\begin{equation}
 q_M(p)=|p|^2I-pp^T+O(|p|^4),\qquad \beta(S_M)=1.
 \label{eqr-reference-beta}
\end{equation}
This normalization was proved by the flat symbol calculation;
ellipticity alone would not imply it. The fixed-outer-radius
classical support estimate and its zero total linear jet also give
\begin{equation}
                    \|[S_M]\|_{\mathcal Q,b}\le C_Sb^2,
 \label{eqr-reference-size}
\end{equation}
uniformly in $M$ and in the subsequent inner radius $b$.
An arbitrary Maxwell-normalized local Wilson reference need not obey
\eqref{eqr-reference-transport} as a full nonlinear functional.

\subsection{The projection telescope}
\label{eqr-projection-section}

Write $\mathcal Q_j$ for the invariant quotient at stage $j$ and set
\begin{equation}
 S_j=[S_{M_j}],\qquad
 P_j=I-S_j\beta_j,\qquad
 \mathcal T_j=\mathcal T_{M_j,R},\qquad
 A_j=P_{j+1}\mathcal T_j.
 \label{eqr-projections}
\end{equation}
The notation $S_j\beta_j$ denotes a rank-one operator.
Because $\beta_j(S_j)=1$, $P_j$ projects onto $\ker\beta_j$.
The established coefficient and reference estimates give
\begin{equation}
 \|\beta_j\|\le C_\beta b^{-2},\qquad
 \|P_j\|\le p:=1+C_\beta C_S.
 \label{eqr-projection-norm}
\end{equation}
The $b^{-2}$ factor has been cancelled by
\eqref{eqr-reference-size}, not omitted.

\begin{lemma}[Cancellation of every intermediate extraction]
\label{eqr-telescope-lemma}
The actual references satisfy
\begin{equation}
 P_{j+1}\mathcal T_jP_j=P_{j+1}\mathcal T_j.
 \label{eqr-intertwining}
\end{equation}
Consequently, for every finite $q$,
\begin{align}
 A_{j+q-1}\cdots A_j
   &=P_{j+q}\mathcal T_{j+q-1}\cdots\mathcal T_j,
       \label{eqr-product}\\
 A_{j+q-1}\cdots A_jP_j
   &=P_{j+q}\mathcal T_{M_j,R^q}P_j.
       \label{eqr-direct-product}
\end{align}
The second identity uses the classical semigroup on the same joint
invariant system.
\end{lemma}
\begin{proof}
For $x\in\mathcal Q_j$, the difference between the two sides of
\eqref{eqr-intertwining} is
$-\beta_j(x)P_{j+1}\mathcal T_jS_j$.
By \eqref{eqr-reference-transport}, this is
$-\beta_j(x)P_{j+1}S_{j+1}=0$.
Repeated application removes the intermediate projections.
The actual scalar semigroup modulo gauge, together with the
universal summable-gauge identity of Section~\ref{csgq-section},
gives \eqref{eqr-direct-product}.
\end{proof}

The proof does not require
$\beta_{j+1}\mathcal T_j=\beta_j$.
Nor does it require equality of sequential and direct carrier arrays.
The common scalar class may be estimated using its separately
justified direct connected representation. If instead
$\mathcal T_jS_j=S_{j+1}+D_j$, the error in
\eqref{eqr-intertwining} is
$-P_{j+1}D_j\beta_j$, potentially costing
$b^{-2}\|D_j\|$. The actual identity makes that error exactly zero.

Use the supported low-kernel and high-jet transport estimates on the
same infinite system. Their explicit common-domain and reserve
hypotheses give
\begin{equation}
 \left\|A_{j+q-1}\cdots A_j\big|_{\ker\beta_j}\right\|
 \le\theta_q:=
 p^2\left[C_{\rm low}(J)R^{-q\alpha}
                     +6C_{\rm high}R^{-q}\right].
 \label{eqr-classical-block}
\end{equation}
The fixed grouping $J$ is chosen before $q$; its finite constant is
not assumed uniform as $J$ grows. Choose a finite $q$ so that
$\theta_q\le1/4$. Then choose $b$ to meet the direct $R^q$
high-jet disk and all classical margins. The constants in
\eqref{eqr-classical-block} do not deteriorate just because $b$
is reduced. Bounded one-step transport gives constants
\begin{equation}
 \|\mathcal T_j\|\le B_T,\qquad
                 \|A_j\|\le B,\qquad B\ge\max(1,pB_T),
 \label{eqr-one-step-linear}
\end{equation}
uniform in $j$. These may be large; no one-step classical
contraction is asserted.

\subsection{The local quantum remainder at fixed coupling}
\label{eqr-quantum-section}

Let $\mathcal F_j$ be the actual limiting local quotient map with
specified coupling $g_j$, retained threshold $T_{{\rm cut},j}$,
and the fixed classical system. Write
\begin{equation}
 \mathcal F_j(x)=\mathcal T_jx+\mathcal G_j(x),\qquad
 \Psi_j(x):=P_{j+1}\mathcal F_j(x)=A_jx+H_j(x),
 \qquad H_j=P_{j+1}\mathcal G_j.
 \label{eqr-quantum-split}
\end{equation}
The remainder includes the pivot and Gaussian determinant terms and
the full base Gaussian logarithm. Those terms are not discarded.
The output of $\Psi_j$ belongs to $\ker\beta_{j+1}$.

Let $\bar\epsilon$ be the representative ceiling from the strict
two-input comparison. The quotient maps are defined on the ball of
radius $\bar\epsilon/4$. Equation
\eqref{qll-quantum-Lipschitz} gives
\begin{equation}
 \|\mathcal G_j(x)-\mathcal G_j(y)\|_{\mathcal Q}
 \le L_{{\rm quant},j}\|x-y\|_{\mathcal Q},\qquad
 L_{{\rm quant},j}
       =C_{\rm quant}\frac{g_jT_{{\rm cut},j}}b.
 \label{eqr-quantum-lipschitz}
\end{equation}
It has no equal-low-jet hypothesis. One can see the domain margin
directly from \eqref{qll-compatible-representatives}: choose
covariant near-minimal representatives for $y$ and $x-y$ and add
them. Both complete inputs then have norm less than
$\bar\epsilon$. Independent near-minimal representatives for $x$
and $y$ would not control their difference.

Choose a uniform number $\nu$ with
\begin{equation}
 pL_{{\rm quant},j}\le\nu,\qquad
 \|H_j(x)-H_j(y)\|\le\nu\|x-y\|.
 \label{eqr-projected-lipschitz}
\end{equation}
All comparisons here hold the coupling, threshold, classical saddle
and cutoff data fixed at each step. The limiting map is well defined
by Theorem~\ref{qll-descent-theorem}; the finite comparison and its
passage to the quotient are provided by
Section~\ref{qll-Lipschitz-section}.

There are two useful forcing estimates. The original common-domain
bound gives
\begin{align}
 \|H_j(0)\|&\le f, \notag\\
 f&\le p\sup_j\left[
 (C_\Delta+C_{\rm det})b+
 2S_1\{e^3\tau(T_{{\rm cut},j})+e^{v_j}v_j\}\right],
 \notag\\
 v_j&\le C_{\rm base}
       (g_jT_{{\rm cut},j}^3+g_jT_{{\rm cut},j}
                              +g_j^2T_{{\rm cut},j}^2).
 \label{eqr-crude-force}
\end{align}
More strongly, fix the outer zero-input construction of
Section~\ref{zif-section} and then restrict to $b\le\rho/2$.
Equations \eqref{zif-force-bound} and
\eqref{zif-extracted-beta} yield
\begin{equation}
 \|\mathcal G_j(0)\|\le C_Fb^2,\qquad
 \|H_j(0)\|\le C_0b^2,\qquad
 C_F=\frac{2H_*}{\rho^2},\quad C_0=pC_F.
 \label{eqr-quadratic-force}
\end{equation}
These are the same actual maps restricted from a common outer
domain, not radius-dependent replacements. The force need not be
$O(g_j^2)$.

\subsection{A finite-block perturbation estimate}
\label{eqr-block-section}

\begin{lemma}[Duhamel bound with a contractive classical block]
\label{eqr-Duhamel-lemma}
Suppose the maps in \eqref{eqr-quantum-split} are evaluated on
their common admissible neighborhood,
\eqref{eqr-classical-block} and \eqref{eqr-one-step-linear} hold,
and \eqref{eqr-projected-lipschitz} holds. For two trajectories
using the same prescribed step couplings and other fixed data,
with their initial difference in $\ker\beta_j$, one has
\begin{align}
 \|\delta_q\|
 &\le\left[
 \theta_q+\nu\sum_{k=0}^{q-1}
                    B^{q-1-k}(B+\nu)^k\right]\|\delta_0\|
       \notag\\
 &=\left[\theta_q+(B+\nu)^q-B^q\right]\|\delta_0\|
       \notag\\
 &\le\left[\theta_q+q\nu(B+\nu)^{q-1}\right]\|\delta_0\|.
 \label{eqr-Duhamel-bound}
\end{align}
The sum formula includes $\nu=0$ by continuity.
If $\|H_i(0)\|\le f$, then any trajectory obeys
\begin{equation}
 \|x_{j+k}\|\le(B+\nu)^k\|x_j\|
                     +f\sum_{l=0}^{k-1}(B+\nu)^l
 \label{eqr-elementary-size}
\end{equation}
as long as the intermediate maps remain admissible.
\end{lemma}
\begin{proof}
The exact variation-of-constants identity is
\begin{align}
 \delta_q={}&A_{j+q-1}\cdots A_j\delta_0 \notag\\
 &+\sum_{k=0}^{q-1}A_{j+q-1}\cdots A_{j+k+1}
   \bigl[H_{j+k}(x_{j+k})-H_{j+k}(y_{j+k})\bigr].
 \label{eqr-Duhamel-identity}
\end{align}
An empty product is the identity. The elementary difference bound
is $\|\delta_k\|\le(B+\nu)^k\|\delta_0\|$.
Apply the improved classical product bound only to the first term
in \eqref{eqr-Duhamel-identity}; bound each shorter product by the
corresponding power of $B$. This proves the first line of
\eqref{eqr-Duhamel-bound}. The difference-of-powers identity gives
the second line, and bounding every sum term by
$(B+\nu)^{q-1}$ gives the third.

For a single step,
$\|\Psi_i(x)\|\le(B+\nu)\|x\|+f$.
Induction proves \eqref{eqr-elementary-size}.
No contraction of a shorter classical product was presumed.
\end{proof}

The leading classical term in \eqref{eqr-Duhamel-bound} has no
growing factor from intermediate projections. The remaining
perturbative sum has finite length $q$, fixed before the coupling
is made small.

\subsection{An invariant quadratic amplitude envelope}
\label{eqr-envelope-section}

\begin{theorem}[Extracted finite-block contraction on the local quotient]
\label{eqr-envelope-theorem}
Use the selected joint infinite homogeneous system, the actual
reference family, the transport estimates
\eqref{eqr-classical-block}--\eqref{eqr-one-step-linear},
the quantum maps of Theorem~\ref{qll-descent-theorem}, and the
fixed-outer-radius force of Theorem~\ref{zif-theorem}.
Choose a finite $q$ with $\theta_q\le1/4$. Put
\begin{equation}
 K=B+1,\qquad L_q=\sum_{l=0}^{q-1}K^l,\qquad
 K_{\rm amp}=\max\{1,2C_0L_q\},\qquad
 C_{\rm int}=K^qK_{\rm amp}+C_0L_q.
 \label{eqr-envelope-constants}
\end{equation}
Choose the common positive $b\le\rho/2$ to satisfy all direct
$R^q$ classical and high-jet domain restrictions and
\begin{equation}
                  C_{\rm int}b^2<\bar\epsilon/4.
 \label{eqr-inner-radius}
\end{equation}
Finally choose the weak-coupling range and thresholds so that the
strict two-input budgets, including
\eqref{zif-inherited-budget}, hold and
\begin{equation}
             \nu\le\min\left\{1,\frac1{4qK^{q-1}}\right\}.
 \label{eqr-nu-choice}
\end{equation}

For each initial stage $j$, every $q$-step extracted block maps
\begin{equation}
 \{x\in\ker\beta_j:\ \|x\|\le r_b\},
                \qquad r_b=K_{\rm amp}b^2,
 \label{eqr-invariant-ball}
\end{equation}
into the corresponding radius-$r_b$ ball at stage $j+q$, with
Lipschitz constant at most $1/2$. Every intervening amplitude is
at most $C_{\rm int}b^2<\bar\epsilon/4$.

For a prescribed countable sequence of couplings in this weak
range, with its fixed common classical local-limit system,
the extracted quotient recursion is therefore defined for all
stages starting in \eqref{eqr-invariant-ball}. At block endpoints
two trajectories using that same sequence contract by $2^{-n}$
after $n$ blocks. These conclusions also bound a density-preserving
running trajectory for as long as it stays in the prescribed
weak-coupling regime.
\end{theorem}
\begin{proof}
Start in \eqref{eqr-invariant-ball}. For $0\le k\le q$,
\eqref{eqr-elementary-size}, $\nu\le1$, and
\eqref{eqr-quadratic-force} give
\begin{equation}
 \|x_{j+k}\|
 \le K^kK_{\rm amp}b^2+C_0b^2\sum_{l=0}^{k-1}K^l
 \le C_{\rm int}b^2<\bar\epsilon/4.
 \label{eqr-intermediate-bound}
\end{equation}
This induction verifies, before the next step is invoked, that
each input lies in the common admissible neighborhood. The
same holds for the trajectory starting at zero and for any
second input in the ball.

Equation \eqref{eqr-Duhamel-bound} and
\eqref{eqr-nu-choice} make the block Lipschitz constant at most
$1/4+1/4=1/2$. Its zero-input output has norm at most
$C_0L_qb^2\le r_b/2$. Therefore its output at any $x$ in
\eqref{eqr-invariant-ball} has norm at most
$\|x\|/2+r_b/2\le r_b$.
The output is extracted at every step by definition.

All estimates are uniform in the initial homogeneous depth.
Repeat the argument block by block. This constructs the
stated infinite sequence of quotient classes for fixed
prescribed step data; it asserts no convergence of a single
nonautonomous trajectory to a fixed point. Applying the same
block Lipschitz bound to two trajectories with identical data
gives the $2^{-n}$ difference estimate.
\end{proof}

The parameter order is noncircular. Choose the primitive data,
the elementary $R$ and grouping $J$ with the strict
Section~\ref{qll-margins-section} margins, then choose the finite
$q$ using \eqref{eqr-classical-block}. The outer radius and
$H_*$ in Section~\ref{zif-section} are fixed before the final
inner radius. Then choose $b$ by
\eqref{eqr-inner-radius} and the direct $R^q$ restrictions.
Only afterward decrease the coupling ceiling. For example,
thresholds equal to a sufficiently large fixed multiple of
$\sqrt{\log(1/g)}$, with a fixed lower threshold, have
$gT(g)/b\to0$, $gT(g)^3\to0$ and $\tau(T(g))\to0$.
Thus the comparison and forcing budgets can be met uniformly
on $0<g\le g_0$. In particular
\[
 \frac{16C_TC_RgT(g)}b\,\bar\epsilon
                   \le\frac{\eta_*}{128e}
\]
is imposed with the stronger two-input budget.
The constants $p,B,\theta_q$ do not grow just because $b$
is reduced.

For comparison, the weaker force \eqref{eqr-crude-force} also
suffices for a fixed-radius envelope. With the same $K,L_q$,
take
\begin{equation}
 r=\frac{\bar\epsilon}{8K^q},\qquad
 f\le\frac{r}{2L_q}.
 \label{eqr-crude-envelope}
\end{equation}
Then intermediate norms are at most
$\bar\epsilon/8+r/2<\bar\epsilon/4$, and the same proof
preserves the radius-$r$ extracted ball. The quadratic
version is stronger because it controls the Maxwell
increment independently of the inner radius.

\subsection{Density bookkeeping and bounded coupling increments}
\label{eqr-coupling-section}

Along the quadratic envelope the unprojected one-step value obeys
\begin{equation}
 \|\mathcal F_j(x_j)\|
 \le \left[B_TC_{\rm int}+C_F
                       +L_{{\rm quant},j}C_{\rm int}\right]b^2.
 \label{eqr-unprojected-size}
\end{equation}
This follows by separating $\mathcal T_jx_j$,
$\mathcal G_j(0)$ and
$\mathcal G_j(x_j)-\mathcal G_j(0)$.
The norm of $\beta$ consequently gives
\begin{equation}
 |\beta_{j+1}(\mathcal F_j(x_j))|
 \le C_\beta\left[B_TC_{\rm int}+C_F
                       +L_{{\rm quant},j}C_{\rm int}\right].
 \label{eqr-beta-increment-bound}
\end{equation}
The factor $b^{-2}$ has again been paid by a proved $b^2$
amplitude. This is an absolute bound, with no sign.

For the actual density, write $\gamma_j=g_j^{-2}$. The
one-step classical factorization is
\[
        \exp[-\gamma_j S_{j+1}+\mathcal F_j(x_j)].
\]
Full-reference extraction gives the exact identity
\begin{equation}
 \mathcal F_j(x_j)
  =\beta_{j+1}(\mathcal F_j(x_j))S_{j+1}+\Psi_j(x_j),
 \qquad
 \gamma_{j+1}
  =\gamma_j-\beta_{j+1}(\mathcal F_j(x_j)).
 \label{eqr-density-running}
\end{equation}
The negative sign is forced by the displayed classical exponent.
A freely prescribed weak coupling sequence need not satisfy
\eqref{eqr-density-running}. The theorem bounds the extracted
amplitudes uniformly over such sequences as operator estimates;
it identifies a continuation of the same density only along
a sequence satisfying this extra update.

Two trajectories that separately update their couplings need
not use the same maps: their coupling and threshold parameters
also differ. Equation \eqref{eqr-quantum-lipschitz} controls only
the inherited-input dependence at fixed coupling. Thus the
$1/2$ block estimate is not a contraction in the joint
coupling--activity variables.

For a fixed finite number of steps,
\eqref{eqr-beta-increment-bound} permits choosing the initial
inverse coupling sufficiently large to retain a positive
weak-coupling range. It does not control an infinite sum of
increments, establish their sign, or prove asymptotic freedom.
It does not prove that a physical running trajectory remains
in the regime needed by Theorem~\ref{eqr-envelope-theorem}.

\subsection{What the quotient envelope does not supply}
\label{eqr-scope-section}

The conclusions concern the invariant infinite quotient of
the actual homogeneous local germ and its specified analytic
continuation. They hold for every admissible input class
on one fixed joint classical system. No uniqueness across
arbitrary classical regulator subsequences is asserted.

Theorem~\ref{qll-varying-theorem} identifies fixed finite
histories only when their actual finite intermediate input
families satisfy the uniform budgets and insertion-domain
hypotheses. A small quotient norm, or a newly localized
infinite representative, does not automatically provide those
bounds for the actual finite winding remainders. For a fixed
finite $q$, crude canonical estimates and an additional radius
reduction may furnish a finite history, but that construction
must be supplied separately. The infinite quotient envelope
is not a regulator-uniform finite RG induction.

If extraction is performed before the local limit, it must
use the actual finite reference $S_{M,N}$ and a specified
finite coefficient functional $\beta_N$ with the required
normalization and convergence. The finite-volume Ward
construction supplies a candidate cutoff-jet functional.
One must not assume that $S_{M,N}$ equals the canonical
periodization of an independently selected infinite reference.
The exact classical identity holds for the actual finite
references; compatibility of projected finite histories
requires the separate coefficient and family statements.

The quotient removes field-independent totals. Physical
vacuum factors remain separately tracked, and their
extensive limits are not supplied by this theorem.
Large-field and mixed histories, global measure consistency,
observables and the continuum limit remain outside this
local amplitude result.

The intermediate projection obstruction is therefore resolved:
the actual classical reference transports exactly, so its
extractions telescope. Combined with the proved local quantum
quotient map and quadratic force, this yields the finite-block
contraction and invariant amplitude envelope above. It does
not constitute a completed Yang--Mills construction.
